% ============ 结语：20 世纪的数论：皇后与仆人 (原书 pp.389-402) ============
\chapter*{结语：20 世纪的数论：皇后与仆人}

$1900$ 年, 希尔伯特（Hilbert）在巴黎举行的第二届国际数学家大会上提出了 $23$ 个著名的数学问题, 其中有 $4$ 个属于代数数论范围. 下面试图以这 $4$ 个问题为线索, 扼要地介绍 $20$ 世纪代数数论的发展.

(1) 希尔伯特第 $12$ 问题是: 对于每个代数数域 $K$, 如何明显构作 $K$ 的最大阿贝尔扩域 $K^{\mathrm{ab}}$？在理论上, $K^{\mathrm{ab}}$ 是 $K$（在 $\C$ 中）的所有有限阿贝尔扩域的合成, 这是 $K$ 的无限扩张. 我们需要对每个代数数域 $K$, 用 $K$ 的某种性质来刻画 $K$ 的所有的有限阿贝尔扩张. 本书第二章中提到, $19$ 世纪末德国人克罗内克（Kronecker）和韦伯（Weber）证明了, 有理数域 $\Q$ 的每个有限阿贝尔扩张 $L$ 都是某个分圆域 $\Q(\zeta_m)$ 的子域, 所以 $\Q$ 的最大阿贝尔扩域 $\Q^{\mathrm{ab}}$ 就是所有分圆域的合成, 即
\[
  \Q^{\mathrm{ab}}=\bigcup_{m\geqslant 3}\Q(\zeta_m)
  =\Q\left(e^{2\pi ia}:a\in\Q\right).
\]
换句话说, $\Q^{\mathrm{ab}}$ 是在 $\Q$ 上添加乘法群 $\C^{\times}$ 的全部有限阶元素（即单位根）而得到的域, 也就是添加周期为 $1$ 的函数 $e^{2\pi ix}$ 在所有有理数点 $x\in\Q$ 的值而得到的域. 希尔伯特第 $12$ 问题是说: 对于任意代数数域 $K$, 是否有类似地方式明显地构作出 $K^{\mathrm{ab}}$？$1920$ 年, 在德国工作的日本数学家高木贞治建立了类域论, 用 $K$ 的``广义''类群来刻画 $K$ 的所有有限阿贝尔扩张. 在这个理论的指导下, 不久对于所有虚二次域 $K=\Q(\sqrt{-d})\ (d>0)$ 明显构作出 $K^{\mathrm{ab}}$, 它是将 $K$ 添加某种双周期函数（椭圆模函数）在特殊点的值而得到的域. 除了有理数域和虚二次域之外, 对于其他数域 $K$ 目前仍不知如何明显构作最大阿贝尔扩域 $K^{\mathrm{ab}}$.

(2) 希尔伯特第 $11$ 问题是问代数数域 $K$ 上多元二次型
\[
  f(x_1,\cdots,x_k)=\sum_{i,j=1}^{k}a_{ij}x_ix_j\qquad(a_{ij}\in K)
\]
都可以表示域 $K$ 的哪些元素？即对于 $K$ 中哪些元素 $\alpha$, 方程 $f(x_1,\cdots,x_k)=\alpha$ 在域 $K$ 中有解？

这个问题于 $1930$ 年由德国数学家哈塞（Hasse, $1898$--$1981$）所解决. 我们在第六章介绍了 $K$ 为有理数情形的哈塞结果. 事实上, 哈塞对任意代数数域 $K$ 证明了同样的结果. $K$ 有无限多个素除子（即 $K$ 上的赋值等价类）. 对于 $K$ 的每个素除子 $P$, 都可以作出局部域 $K_P$, 这些局部域都是 $K$ 的扩域. 所以对每个 $\alpha\in K^{\times}$, 如果 $f(x_1,\cdots,x_k)=\alpha$ 在 $K$ 中有解, 那么它也是每个局部域 $K_P$ 中的解. 反过来, 哈塞证明了: 若对每个素除子 $P$（包括 $K$ 的无限素除子）, 方程 $f(x_1,\cdots,x_n)=\alpha$ 在 $K_P$ 中均有解, 则此方程一定在 $K$ 中有解. 类似地, 方程 $f(x_1,\cdots,x_n)=0$ 在 $K$ 中有非零解, 当且仅当此方程在每个局部域 $K_P$ 中有非零解. 哈塞把问题化为局部域上的事情（局部--整体原则）. 进而, 对每个具体的二次型方程 $f(x_1,\cdots,x_n)=\alpha$, 它在几乎所有（即除有限个之外）的局部域中均有解, 从而只需考查它在有限个局部域 $K_P$ 中是否有解即可. 而对每个局部域 $K_P$, 方程 $f(x_1,\cdots,x_n)=\alpha$ 在 $K_P$ 中是否有解则有相当方便的判别方法（希尔伯特符号）. 这就至少在理论上相当漂亮地回答了希尔伯特第 $11$ 问题. 例如当 $n\geqslant 5$ 时, 用希尔伯特符号可以证明, $f(x_1,\cdots,x_n)=\alpha$（对每个 $\alpha\in K^{\times}$）在 $K_P$（$P$ 为 $K$ 的有限素除子, 即为 $O_K$ 的素理想）中一定有解, 从而只要 $f(x_1,\cdots,x_n)=\alpha$ 对于无限素除子 $P$ 的局部域 $K_P$（它为 $\R$ 或 $\C$）中有解, 则在 $K$ 中必有解.

哈塞的工作之后, 局部域理论和 $p$-adic 分析逐渐成为代数数论的主要方法和工具. 这种方法的另一个深刻影响是把数论和几何联系在一起. 设 $f(x,y)$ 是实系数多项式, 则 $C\colon f(x,y)=0$ 是实坐标平面的一条曲线. 域 $\R(x,y)$ 叫作曲线 $C$ 的函数域. 由于 $y$ 满足 $f(x,y)=0$, 即 $\R(x,y)$ 是有理函数域 $\R(x)$ 的代数扩张. 如果把实数域 $\R$ 改成有限域 $\F_q$, 而 $f(x,y)\in\F_q[x,y]$, 则
\[
  C\colon f(x,y)=0
\]
是 $\F_q$ 上的一条代数曲线, 它的函数域 $\F_q(C)$ 为 $\F_q(x,y)$. 德国数学家阿廷（Artin）发现这种函数域 $K=\F_q(C)$ 有和代数数域同样的赋值理论, 从而也有一系列局部域 $K_P$（注意: $K$ 的素除子 $P$ 均是非阿的）. 由此得到函数域 $K=\F_q(x,y)$ 的许多代数性质与数域相像, 并且可用来研究有限域上代数曲线的几何性质. 比如说, 上述关于二次型方程的哈塞局部--整体原则对于函数域也是对的. 此后人们把数域和（有限域上代数曲线的）函数域统称为整体域, 而后来的数论常常统一考虑整体域上的问题. 几何与数论开始结合在一起, $50$ 年代后逐渐发展成一个重要的方向, 叫作算术代数几何. 它的第 $1$ 项重要成果是在希尔伯特第 $8$ 问题上取得的.

(3) 希尔伯特第 $8$ 问题是问: 在整数环 $\Z$ 和有理数域范围内的许多数论问题（如黎曼猜想, 素数分布, 哥德巴赫猜想等）如何推广到任意代数数域上去？本世纪初, 沿用已有的解析工具, 素数分布的结果平行地推广成代数数域的整数环素理想分布结果. 黎曼猜想也可推广成``广义''黎曼猜想, 即对每个代数数域 $K$, 戴德金 zeta 函数 $\zeta_K(s)$ 的非平凡零点的实数部分猜想均为 $\frac{1}{2}$. 到目前为止, 对于任何代数数域 $K$, 这个广义黎曼猜想都没有证明. 可是对于函数域的情形, 黎曼猜想的一个模拟却在 $1948$ 年被证明!

对于有限域 $\F_q$ 上代数曲线 $C\colon f(x,y)=0$, 我们有函数域 $K=\F_q(C)=\F_q(x,y)$. 可以像数域那样定义域 $K$ 的 zeta 函数
\[
  \zeta_K(s)=\prod_{\mathfrak{p}}(1-N(\mathfrak{p})^{-s})^{-1}
  \qquad(s=\sigma+it,\ \sigma>1),
\]
其中 $\mathfrak{p}$ 过域 $K$ 的所有素除子. 法国数学家韦伊（Andr\'e Weil）运用几何上的黎曼--罗赫定理, 证明了这个函数是变量 $U=q^{-s}$ 的有理函数: 将 $\zeta_K(s)$ 写成 $Z_K(U)$ 之后,
\[
  Z_K(U)=\frac{L(U)}{(1-U)(1-qU)},
\]
其中
\[
  L(U)=1+c_1U+\cdots+c_{2g-1}U^{2g}+q^gU^{2g}\in\Z[U],
\]
这里 $g=g(K)=g(C)$ 是曲线 $C$ 的亏格（genus）, 这是 $C$ 的双有理变换的一个重要的几何不变量, 也叫 $K$ 的亏格（是非负整数）. 于是 $Z_K(U)$ 是关于 $U$ 的亚纯函数. 韦伊还证明了函数方程
\[
  Z_K(U)=(\sqrt{q}U)^{2g}Z_K\left(\frac{1}{qU}\right)^{-1},
\]
这也相当于 $L(U)=(\sqrt{q}U)^{2g}L\left(\frac{1}{qU}\right)$. 若令
\[
  L(U)=\prod_{i=1}^{2g}(1-\alpha_iU)\qquad(\alpha_i\in\C),
\]
即 $\alpha_i^{-1}\ (1\leqslant i\leqslant 2g)$ 是 $L(U)$（和 $Z_K(U)$）的全部零点, 由函数方程可知（适当安排零点的次序）$|\alpha_{2g-i+1}\alpha_i|=q\ (1\leqslant i\leqslant g)$. $1941$ 年, 韦伊根据一些例子, 提出如下的猜想:
\[
  |\alpha_i|=\sqrt{q}\qquad(1\leqslant i\leqslant 2g).
\]
考虑到 $\alpha_i^{-1}\ (1\leqslant i\leqslant 2g)$ 是 $Z_K(U)$ 的零点以及 $U=q^{-s}$, 可知韦伊猜想相当于说 $\zeta_K(s)$ 的零点满足 $q^{-s}=\alpha_i^{-1}\ (1\leqslant i\leqslant g)$, 即 $s$ 的实数部分均为 $\frac{1}{2}$. 所以韦伊猜想是黎曼猜想在函数域情形的一种模拟. $1948$ 年, 韦伊本人证明了这个猜想. 为了证明此猜想, 韦伊创造了代数几何的许多新概念, 把代数几何推进到一个新的阶段. 到了 $1979$ 年才出现了只用黎曼--罗赫定理来证明韦伊定理的简化证明.

韦伊的结果有重要的理论价值, 可以计算和估计一些重要的指数和, 对有限域上曲线的点数给出很好的估计, 而且还可用于通信和计算机科学中. 设 $C\colon f(x,y)=0$ 是 $\F_q$ 上的一条代数曲线. 对每个 $n\geqslant 1$, 以 $N_n=N_n(\mathcal{C})$ 表示此曲线在 $\F_{q^n}$ 中的射影点数（即方程 $f(x,y)=0$ 在 $\F_{q^n}$ 中的解数, 包含其上的``无穷远点''）, 我们可以构作另一个 zeta 函数
\[
  Z_{\mathcal{C}}(U)=\exp\sum_{n\geqslant 1}\frac{N_n}{n}U^n.
\]
可以证明, 当 $\mathcal{C}$ 是绝对不可约的光滑曲线时, 这个 zeta 函数等于前面定义的 $Z_K(U)$（其中 $K=\F_q(x,y)$ 是曲线 $\mathcal{C}$ 的函数域）. 由此得到
\[
  N_n=q^n+1-\sum_{i=1}^{2g}\alpha_i^n.
\]
再由韦伊定理便得到
\[
  |N_n-(q^n+1)|\leqslant 2gq^{n/2}.
\]
事实上, 韦伊对于有限域上高维代数簇给出类似的猜想. 在 $20$ 世纪后半期代数几何建立了概型理论和各种上同调理论之后, 高维韦伊猜想才于 $1973$ 年由德林（Deligne）证明.

(4) 除了类域论和算术代数几何之外, $20$ 世纪的前半期还产生了数论的另外两个重要分支: 模形式理论和椭圆曲线的算术理论.

以 $H=\{z=\sigma+it\in\C\mid t>0\}$ 表示上半平面, 群
\[
  \Gamma=\left\{g=\begin{pmatrix}a&b\\ c&d\end{pmatrix}\mid a,b,c,d\in\Z,\ ad-bc=1\right\}
\]
在 $H$ 上的作用为
\[
  g(z)=\frac{az+b}{cz+d}\qquad\left(g=\begin{pmatrix}a&b\\ c&d\end{pmatrix}\in\Gamma\right).
\]
易知 $\Gamma$ 中元素 $g$ 把上半平面 $H$ 变成自身. 定义于 $H$ 上的复变解析函数 $f(z)$ 叫作对于 $\Gamma$ 的某个子群 $G$ 的权 $k$ 模形式, 是指对每个 $g=\begin{pmatrix}a&b\\ c&d\end{pmatrix}\in G$,
\[
  f\left(\frac{az+b}{cz+d}\right)=(cz+d)^kf(z).
\]
全体对于群 $G$ 的权 $k$ 模形式构成的集合 $M_k(G)$ 是复数域 $\C$ 上的有限维向量空间. 模形式理论的最基本问题是计算各种模形式空间的维数和构作一组基. 在 $20$ 世纪 $20$ 年代, 德国数学家 Hecke 得到一个重要结果. 对于 $\Gamma$ 的某些重要的子群 $G$, 存在正整数 $M$, 使得 $\begin{pmatrix}1&M\\ 0&1\end{pmatrix}\in G$. 于是模形式空间 $M_k(G)$ 中的每个模形式都满足 $f(z+M)=f(z)$, 从而有 Fourier 展开 $f(z)=\sum_{n=0}^{\infty}a_ne^{2\pi iz/M}$.

Hecke 证明了: 空间 $M_k(G)$ 存在一组特殊的基, 其每个基元素 $f(z)=\sum_{n=0}^{\infty}a_ne^{2\pi iz/M}$ 所对应的 $L$-函数 $L_f(s)=\sum_{n=1}^{\infty}a_nn^{-s}$ 有欧拉无穷乘积展开
\[
  L_f(s)=\prod_p(1-a_pp^{-s}+p^{k-1-2s})^{-1},
\]
其中 $p$ 过所有素数. 并且有形如 $L_f(s)=F(s)L_f(k-s)$ 的函数方程, $F(s)$ 是一个简单的函数. 具有这种性质的模形式叫作 Hecke 形式.

模形式理论的最早起源仍是希尔伯特第 $11$ 问题, 即研究二次型方程 $f(x_1,\cdots,x_k)=\sum_{i,j=1}^{k}a_{ij}x_ix_j=\alpha$ 的解, 不过这时 $a_{ij}$ 和 $\alpha$ 均属于某个代数数域 $K$ 的整数环 $O_K$, 而问题是求方程在 $O_K$ 中的解. 这是比求域 $K$ 中的解更困难的问题. 比如高斯研究二元二次型 $ax^2+bxy+cy^2=n\ (a,b,c,n\in\Z)$ 的整数解, 就是 $K=\Q$（$O_K=\Z$）和 $n=2$ 的情形. 如果 $f(x_1,\cdots,x_k)$ 是 $\Z$ 上的正定二次型, 则对每个正整数 $n$, 方程 $f(x_1,\cdots,x_k)=n$ 均有有限多个整数解. 以 $a_n$ 表示其解数, 则对于许多二次型 $f$, 函数 $F(z)=\sum_{n=1}^{\infty}a_ne^{2\pi iz}$ 都是对某个群 $G$ 的权 $k$ 模形式. 假如我们知道模形式空间 $M_k(G)$ 的一组基, 通过有限个 $a_n$ 的值我们可以把 $F(z)$ 表成这组基的线性组合, 于是便可得到方程 $f(x_1,\cdots,x_k)=n$ 的解数 $a_n$ 的一般公式.

现在介绍另一个数论分支: 椭圆曲线的算术理论. 设 $K$ 是代数数域, $C\colon f(x,y)=0$ 是 $K$ 上一条代数曲线. 在代数几何中, 我们研究曲线的双有理分类. 其中一个重要的双有理不变量叫作曲线的亏格 $g=g(C)$. 在代数几何中, 亏格不同的曲线属于不同的双有理等价类. 但是从数论的角度, $K$ 上所有曲线一共分成 $3$ 大类:

（一）$g(C)=0$ 的曲线叫作有理曲线, 这类曲线包括直线和所有二次曲线, 这种曲线上一定有无限多个 $K$-点（即坐标属于 $K$ 的点）.

（二）$g(C)\geqslant 2$. $1923$ 年, 英国数学家莫德尔（Mordell）猜想: 这种曲线只有有限多个 $K$-点. 这个猜想于 $1983$ 年由德国的法廷斯（Faltings）所证明.

（三）$g(C)=1$, 这种曲线叫作椭圆曲线. 它的典型方程为 $E\colon y^2=x^3+ax+b\ (a,b\in K)$, 并且 $4a^3+27b^2\neq 0$. 这种曲线上的 $K$-点个数可以有限也可以无限. 以 $E(K)$ 表示椭圆曲线 $E$ 的所有 $K$-点所成的集合. 可以定义一种运算使得 $E(K)$ 为交换群. 莫德尔证明了 $E(K)$ 是有限生成的交换群, 所以
\[
  E(K)=E(K)_f+E(K)_t\qquad(\text{直和}),
\]
其中 $E(K)_t$ 是有限（交换）群, 而 $E(K)_f$ 是 $E(K)$ 的自由部分, 即 $E(K)_f$ 同构于加法群 $\Z^r$, 其中 $r=r(E)$ 叫作 $K$ 上椭圆曲线 $E$ 的秩. 当 $r=0$ 时 $E(K)$ 有限, 即方程 $y^2=x^3+ax+b$ 在 $K$ 中只有有限多解; 而 $r\geqslant 1$ 时 $E(K)$ 无限, 即方程在 $K$ 中有无限多解. 椭圆曲线算术理论的基本问题是决定 $E(K)_t$, 秩 $r(E)$ 以及给出 $E(K)_f$ 的一组基. 对每条椭圆曲线 $E$, 决定 $E(K)_t$ 比较容易, 而后两个问题则比较困难. 关于群 $E(K)$ 的结构有两个基本猜想:

(i) 对每个固定的代数数域 $K$, $K$ 上所有椭圆曲线 $E$, $E(K)_t$ 只有有限多可能的群结构. $1981$ 年, 美国哈佛大学的梅祖尔（B. Mazur）对于 $K=\Q$ 的情形证明了这个猜想: $E(\Q)_t$ 只有以下 $15$ 种可能
\[
  \Z/n\Z\ (n=1,2,3,\cdots,9,10,12),\qquad
  \Z/2n\Z\oplus\Z/2\Z\ (n=1,2,3,4).
\]
对于任意代数数域 $K$, 这个猜想最终于 $1995$ 年被证明.

(ii) 对于每个固定的代数数域 $K$, $K$ 上所有椭圆曲线 $E$ 的秩 $r(E)$ 是没有上界的. 这个猜想至今没有解决.

为了研究群 $E(K)$ 的结构, 人们使用了许多代数工具（局部域上的椭圆曲线, Tate 模和其上的表示, Selmer 群和 Shafarevich 群...）, 也采用解析数论的手段. 就像用黎曼 zeta 函数 $\zeta(s)$ 和狄利克雷 $L$-函数 $L(s,\chi)$ 来反映素数性质, 用 Dedekind zeta 函数 $\zeta_K(s)$ 来反映代数数域 $K$ 的算术特性, 用 $\zeta_C(s)$ 来反映有限域上曲线 $C$ 的算术性质那样, 人们也希望寻求一个解析函数 $L_E(s)$ 来反映椭圆曲线 $E$ 的算术性质. $1960$ 年前后, 英国数学家基于数值计算, 构作了这样的函数.

我们以 $K=\Q$ 为例. 设 $E\colon y^2=x^3+ax+b\ (a,b\in\Z)$ 是椭圆曲线, $4a^3+27b^2\neq 0$, 则对每个素数 $p$, 当 $p\nmid 4a^3+27b^2$ 时, $E$ 也是 $\F_p$ 上的一条椭圆曲线（即亏格为 $1$）. 所以将 $E$ 看成 $\F_p$ 的椭圆曲线（记成 $E_p$）时, 它的韦伊 zeta 函数为
\[
  \zeta_{E_p}(s)=\frac{F_p(p^{-s})}{(1-p^{1-s})(1-p^{-s})},
\]
其中（$U=p^{-s}$）
\[
  F_p(U)=1-a_pU+pU^2=(1-\alpha_pU)(1-\overline{\alpha}_pU)\qquad(a_p\in\Z).
\]
以 $N_p$ 表示椭圆曲线 $E_p$ 在有限域 $\F_p$ 中的点数, 则 $N_p=p+1-(\alpha_p+\overline{\alpha}_p)=p+1-a_p$, 并且由韦伊定理知 $|a_p|\leqslant|\alpha_p|+|\overline{\alpha}_p|\leqslant 2\sqrt{p}$. $1960$ 年, 英国数学家 Birch 和 Swinnerton-Dyer 两人构作了椭圆曲线 $E$ 的 $L$-函数
\[
  L_E(s)=\prod_p\frac{1}{F_p(p^{-s})}
  =\prod_p\frac{1}{1-a_pp^{-s}+p^{1-2s}}\qquad(s=\sigma+it).
\]
由 $|a_p|\leqslant 2\sqrt{p}$ 可知, 当 $\sigma>\frac{3}{2}$ 时, 这个无穷乘积是收敛的, 即上式定义出 $L_E(s)$ 为右半平面 $\sigma>3/2$ 中的解析函数. 他们猜想:

(BSD1) $L_E(s)$ 可以解析开拓成整个复平面上的解析函数, 并且有形如 $L_E(s)=\varphi(s)L_E(2-s)$ 的函数方程.

函数 $L_E(s)$ 如何反映椭圆曲线 $E$ 的点群 $E(\Q)$ 的性质？上述无穷乘积在 $s=1$ 处不一定收敛, 但是可以形式地表示成
\[
  L_E(1)\sim\prod_p\frac{1}{1-a_pp^{-1}+p^{-1}}
  =\prod_p\frac{p}{p-a_p+1}=\prod_p\frac{p}{N_p}.
\]
上述两位英国数学家的想法是: 如果椭圆曲线 $E_p$ 在有限域 $\F_p$ 中的点数愈多（即 $N_p$ 愈大）, 则相信椭圆曲线 $E$ 在 $\Q$ 中的点数 $|E(\Q)|$ 愈多. 当 $N_p$ 均足够大, 使得 $\prod_p\frac{p}{N_p}=0$, 即 $s=1$ 为 $L_E(s)$ 的零点, 则许多例子表明 $|E(\Q)|=+\infty$（即 $E(\Q)$ 的秩 $r(E)\geqslant 1$）. 更进一步, 他们猜想

(BSD2) $L_E(s)$ 在 $s=1$ 处零点的阶等于群 $E(\Q)$ 的秩 $r(E)$.

当 $L_E(1)\neq 0$ 时, 他们还猜想 $L_E(1)$ 的值和有限交换群 $E(\Q)_t$ 的算术性质有确定的关系. 上述猜想自从 J. Coates 和 A. Wiles $1977$ 年的重要工作以来, 取得很大的进展, 但至今仍未完全解决.

将定义椭圆曲线 $E$ 的 $L$-函数 $L_E(s)$ 的无穷乘积和前面所述 Hecke 模形式（权 $k=2$）对应的 $L$-函数无穷乘积表达式相比, 不难看出它们非常相像. 在 $1950$--$1970$ 年期间, 日本数学家谷山（Taniyama）, 志村（Shimura）和法国数学家韦伊（A. Weil）从不同的角度提出了如下的猜想

\emph{谷山--志村--韦伊猜想:} $\Q$ 上每个椭圆曲线 $E$ 的 $L$-函数 $L_E(s)$ 都是对某个群的权 $2$ Hecke 模形式的 $L$-函数.

这个猜想还有其他等价的表达方式. 在几何上它可以说成是: 每条椭圆曲线都是模曲线在某种映射之下的象. 用表示论的语言则可说成: 由椭圆曲线的 Tate 模构作的伽罗瓦群表示一定是 modular 表示.

以上介绍的 BSD 猜想和谷山--志村--韦伊猜想, 是椭圆曲线的两个重要的猜想. 如果后一个猜想成立, 由于 Hecke 模形式的 $L$-函数已经证明有函数方程并且可以解析开拓, 所以每个椭圆曲线的 $L$-函数也都如此, 即猜想 (BSD1) 成立. 长期以来, 人们一直认为谷山--志村--韦伊猜想是非常困难的. 但是在 $1986$ 年, 德国人 Frey 令人惊奇地发现, 由谷山--志村--韦伊猜想可以推出费马猜想. 自那时起, 怀尔斯经过 $8$ 年的艰苦努力, 对于所谓半稳定（Semistable）椭圆曲线证明了谷山--志村--韦伊猜想正确, 由此就最终证明了具有 $350$ 多年历史的费马猜想. 到了 $1998$ 年, 除了少数情形, 谷山--志村--韦伊猜想几乎完全被证明.

以上我们主要介绍了 $1900$--$1970$ 年近代代数数论的发展轮廓. 这个时期代数数论的研究引入了新的方法和工具（局部--整体原则, 解析方法和几何方法）, 建立了许多新的理论和分支（类域论, 算术代数几何, 模形式理论, 椭圆曲线的算术理论）, 取得了一系列重要成果. 在这段时期, 分圆域的理论也取得重要的发展, 其中最主要的进展是采用 $p$-adic 分析工具引入 $p$-adic zeta 函数和岩泽（Iwasawa）理论, 这些方法和理论后来用于研究椭圆曲线等其他数论问题中.

(5) $1967$ 年, 韦伊出版了《基础数论》（Basic Number Theory）一书. 这本书的内容是讲类域论, 但是采用了全新的讲法, 用单代数和群表示论为工具, 统一讲述整体域（即数域和函数域）, 使用了 Adele 环 $\mathbb{A}_K$ 和 Idele 群 $\mathbb{I}_K$ 的语言, 显示出作者对 $20$ 世纪前 $70$ 年数论发展的深刻理解, 这本书对于后来 $30$ 年数论发展产生很大的影响. 同一年, 美国数学家朗兰兹（Langlands）提出了一系列原则性的重要猜想. 他的最基本看法是用无限群的表示理论研究数论. 可以把它看作是现代的代数数论发展的新时期. 朗兰兹猜想被人称之为一种哲学或纲领. 它的起源则可归根于希尔伯特第 $9$ 问题.

希尔伯特第 $9$ 问题是问: 高斯的二次互反律如何推广到任意代数数域上？

高斯的二次互反律是很初等的. 设 $p$ 是奇素数, $a$ 是与 $p$ 互素的整数. 如果 $a$ 是模 $p$ 的二次剩余（即存在整数 $b$, 使得 $a\equiv b^2\ (\mathrm{mod}\ p)$）, 则记 $\left(\frac{a}{p}\right)=1$, 否则记 $\left(\frac{a}{p}\right)=-1$. 二次互反律是说, 对于任意两个不同的奇素数 $p$ 和 $q$,
\[
  \left(\frac{p}{q}\right)\left(\frac{q}{p}\right)
  =(-1)^{\frac{p-1}{2}\cdot\frac{q-1}{2}}.
\]
高斯给出它的 $6$ 个证明, 并且用二次域中素数的分解法则来解释它, 所以它也可以看成是二次域的互反律: 二次域 $K=\Q(\sqrt{d})$ 中素数 $p$ 的分解模式只依赖于 $p$ 所在的模 $|d(K)|$（$K$ 的判别式）或 $4d(K)$ 的同余类. 或者说: 多项式 $x^2-d$ 分解模式只依赖于 $p$ 所在的模 $|d(K)|$ 或 $4d(K)$ 的同余类. 于是我们可以问: 对于任意整系数不可约多项式 $f(x)\in\Z[x]$, 如何用 $f(x)$ 的自身特性（例如用 $f(x)$ 的系数, 或用 $f(x)$ 的分裂域 $K$）来描述对每个素数 $p$, $f(x)\ (\mathrm{mod}\ p)$ 的分解模式？本书第二章的末尾我们讲到, 如果 $f(x)$ 是 Abel 多项式, 即若 $\alpha$ 是 $f(x)$ 的一个根, 而 $K=\Q(\alpha)$ 为 Abel 数域, 令 $m$ 是域 $K$ 的导子（即 $\Q(\zeta_m)$ 是包含 $K$ 的最小分圆域）, 则 $f(x)\ (\mathrm{mod}\ p)$ 的分解模式只依赖于 $p$ 所在的模 $m$ 同余类. 这可以看成是 Abel 域的互反律. 在 $1920$ 年类域论产生之后, 德国数学家阿廷（E. Artin）用类域论来表达任意数域阿贝尔扩张 $E/K$ 的互反律, 即用 Abel 群 $\Gal(E/K)$ 来描述 $O_K$ 的素理想在 $E$ 中如何分解. 对于一般情形, 即 $f(x)$ 的分裂域为 $K$, 而 $K/\Q$ 的伽罗瓦群不是交换群, 如何刻画 $f(x)\ (\mathrm{mod}\ p)$ 的分解模式？这种情形的复杂之处, 在于群 $\Gal(K/\Q)$ 的不可约复表示不全是一次的. 对于数域的任意伽罗瓦扩张 $E/K$, 和伽罗瓦群 $G=\Gal(E/K)$ 的每个 $n$ 次不可约复表示 $\rho\colon G\to GL_n(\C)$, 阿廷均构作了一个 $L$-函数 $L(E/K,\rho,s)$. 利用有限群表示理论, 阿廷证明了这些 $L$-函数的一系列性质, 并且对它们的解析开拓和解析特性提出一系列猜想. 特别地, 若 $\rho$ 是平凡表示, 则对应的阿廷 $L$-函数就是戴德金 zeta 函数 $\zeta_K(s)$. 而阿廷猜想 $L(E/K,\rho,s)$ 一定是解析的（没有极点）. 利用类域论可以证明这个猜想对于 $L/K$ 是 Abel 扩张时是对的, 对于 $L/K$ 是非阿贝扩张的情形, 目前只对少数情形证明了阿廷的这个猜想. 这些工作表明群表示理论在研究数论问题时开始发挥作用. 朗兰兹的第 $1$ 个猜想就是对阿廷 $L$-函数作出的论断, 他使用了无限群的表示理论. 首先, 他和他的同事将模形式理论以无限群表示论的形式加以改造并作了极大的推广, 吸取了韦伊的思想, 把模形式推广成群 $GL_n(\mathbb{A}_K)$（$\mathbb{A}_K$ 是 $K$ 的 Adele 环）上自守表示和自守形式的深刻理论. 对于群 $GL_n(\mathbb{A}_K)$ 的每个``尖点''自守表示 $\pi$, 朗兰兹构作了一个 $L$-函数 $L(s,\pi)$, 并且证明了当 $\pi$ 为尖点表示时, $L(s,\pi)$ 可延拓成整个复平面上的全纯函数, 有欧拉无穷乘积展开, 并且有函数方程把 $L(s,\pi)$ 和 $L(1-s,\tilde{\pi})$ 联系起来, 其中 $\tilde{\pi}$ 是由 $\pi$ 给出的另一个自守表示.

\emph{朗兰兹第 $1$ 猜想:} 设 $E/K$ 是数域的伽罗瓦扩张, $\rho$ 是伽罗瓦群 $\Gal(E/K)$ 的一个 $n$ 维不可约复表示, 则必存在群 $GL_n(\mathbb{A}_K)$ 的某个尖点表示 $\pi$, 使得阿廷 $L$-函数 $L(E/K,\rho,s)$ 等于 $L(s,\pi)$.

由于 $L(s,\pi)$ 已证明有解析延拓、函数方程等性质, 所以若上述猜想成立, 则阿廷 $L$-函数也有这些性质. 特别地, 当 $\rho\neq 1$ 时, 阿廷 $L$-函数 $L(E/K,\rho,s)$ 是全纯的. 也就是说, 朗兰兹第 $1$ 猜想可以推出阿廷猜想. 当 $E/K$ 是 Abel 扩张时, $\rho$ 是一维表示, 利用类域论已经证明了朗兰兹第一猜想在 $n=1$ 的情形是对的, 而对于一般的 $n$, 朗兰兹这个猜想可以看成是类域论的高维推广. 目前只对少数 $n\geqslant 2$ 维不可约表示朗兰兹的上述猜想被证明是正确的. 此后, 朗兰兹又提出一系列范围更广的猜想. 这些猜想所体现的精神是: 对于数学上的各种研究对象 $X$（例如数论中整体域上的伽罗瓦扩张, 代数几何中的代数簇, 微分几何中的黎曼流形, 动力系统中粒子团, 或者离散数学中的正规图）, 人们试图寻求一个解析函数, 使得此函数的解析特性（零点和极点, 函数方程, 展开系数, 在特殊点的取值...）能够充分反映该数学对象 $X$ 的算术、代数、几何或者组合特性. 朗兰兹认为, 我们一定可以找到一个适当的环 $R(X)$（与 $X$ 有关）, 使得这样的解析函数是 $n$ 维一般线性群 $GL_n(R(X))$ 的某种自守表示 $\pi$ 的 $L$-函数 $L(s,\pi)$.

$20$ 世纪前 $70$ 年, 数论取得许多成就, 而后 $30$ 年更是数论的大丰收时期. 它最集中表现为 $3$ 个菲尔兹奖和两个沃尔夫奖. $1978$ 年, Deligne 由于证明关于有限域上代数簇 zeta 函数的高维韦伊猜想而获菲尔兹奖, 这项工作表明 $1950$ 年以来发展的概型和上同调理论在算术代数几何中发挥巨大的威力. $1986$ 年 Faltings 证明了 Mordell 猜想（亏格大于 $1$ 的代数曲线在整体域中只有有限多解）, 把算术几何的研究又推进到一个新阶段. $1990$ 年 Drinfeld 证明了函数域上二维局部朗兰兹猜想. $1996$ 年, 沃尔夫奖颁给两位数学家, 一位是证明费马猜想的 A. Wiles, 另一位就是朗兰兹.

(6) $20$ 世纪的后 $30$ 年, 数论研究获得了一系列重大的理论成就, 与此同时, 数论在实际领域得到了许多应用. 其中一个重要原因是由于数字计算机的发展和数字通信技术的飞速进步. 计算机科学和数字通信技术为数学提出许多新的研究课题, 这些课题所使用的数学工具不再是传统的傅里叶分析, 而是代数、数论和组合手段. 数论的深刻理论得到应用, 使工程和技术领域发生重大变革. 这里我们只举数字通信中的两个例子.

（一）\emph{代数几何码} 基于通信可靠性的考虑, 从 $20$ 世纪 $50$ 年代人们努力寻找和构作好的纠错码. 利用数论和有限域上代数工具, 建立了线性码（用线性代数）、循环码（用多项式环的理想论）, 最后找到纠错性能良好的 BCH 码. 由于发明了好的译码算法和方便的技术实现手段（移位寄存器）, 从 $60$ 年代以来, BCH 码在数字通信中得到普遍应用. $70$ 年代, 前苏联数学家 Goppa 以黎曼--罗赫定理为工具, 利用有限域上的代数曲线构作了新型纠错码: 代数几何码. $1982$ 年, 欧洲 $3$ 位算术几何学家使用模形式理论的深刻结果, 利用模曲线构作出纠错性能优于 BCH 码的一种代数几何码. 此后十多年来, 代数几何码的译码算法研究取得很大进步, 纠错码的研制和技术使用今后将会更加依赖数论和算术几何的现代理论成果.

（二）\emph{公开密钥体制} 数字通信传输的信息容易被第三者截收和伪造. 随着社会的信息化, 通信的安全性不仅限于政治和军事领域, 已成为电子商务等经济领域和日常生活的重要课题. 通信的网络化产生了大量密钥的保存、分配和管理等诸多问题. 保密通信的传统要求和手段（签名、身份认证、防伪、仲裁等等）在数字通信中如何实现？如何科学评价一个通信系统的安全性？这些实际需要为数学家提出一系列富于挑战性的课题, 促进人们寻找新型的保密通信体制. $1976$ 年, 美国两位年青计算机学家 Diffie 和 Hellman 提出了公开密钥体制, 它解决了大量密钥保存问题和数字签名问题, 受到人们广泛注意. 一时间人们提出了实现这种体制的各种具体方案, 其中有许多方案又相继被人们破掉（即给出技术上可实现的破译算法）. 目前被人们认为可靠的并且已被广泛应用的主要有两种方案, 一种是由美国麻省理工学院 $3$ 位年青数论和计算机学者 Rivest, Shamir 和 Adleman 于 $1977$ 年提出的 RSA 算法, 另一种是离散对数算法. 公开密钥体制被认为是现代保密通信的一场革命. RSA 算法和离散对数算法分别是基于如下的思想: 大整数的因子分解和求离散对数是非常困难的. 所谓离散对数即指对于大素数 $p$, 整数 $a$ 和 $b$, 求整数 $x$ 使得 $a^x\equiv b\ (\mathrm{mod}\ p)$. 使用目前最快的计算机（硬件）和最好的算法（软件）, 分解 $100$ 位的整数需要几分钟, 而分解大约 $200$ 位的整数则需要万年以上. 对于离散对数问题, 目前在素数超过 $150$ 位时也是实际不可能计算的. 但是, 大数分解和求离散对数是否真的很难？是否有更好的算法？从 $70$ 年代以来, 许多优秀的数论和计算机专家热衷于大数分解和离散对数算法的研究, 发明了二次筛法、椭圆曲线算法、数域筛法、函数域筛法等一系列改进的算法, 这些算法使用了现代数论的深刻理论结果. 数论和计算机相结合, 形成了计算数论这个新分支, 从事两种类型的工作. 一种是: 借助于计算机进行数论的理论研究（例如计算数域的理想类数, 寻求基本单位组）, 进行数论的实验, 对于理论研究有所启示和验证. 另一种是数论在计算机科学、通信工程等各种实际领域的应用. 目前, 国际上已经形成了相当强大的计算数论研究队伍和实验室. 几千年前源于古代人类实践的数论, 在经历了辉煌的理论探究里程之后, 正在应用领域发挥着巨大的威力. 数论正在同时扮演着皇后和仆人的角色.
